\documentclass[border=0pt]{standalone}

\usepackage{tikz}        % drawing macros
  \usetikzlibrary{math}
  \usetikzlibrary{calc}  % calculations
  \usetikzlibrary{arrows}  % arrowheads
  \usetikzlibrary{arrows.meta}  % arrowheads
  \usetikzlibrary{decorations.markings}
\usepackage{tikz-3dplot} % drawing macros 3d
\usepackage{pgfplots}
%\pgfplotsset{compat=1.18}

\begin{document}

%\usepackage{tikz}
%\usepackage{tikz-3dplot}
%
\tdplotsetmaincoords{ 75}{ 50}% x-rot,z-rot
\begin{tikzpicture}[scale=2.5,tdplot_main_coords]
  \begin{scope}[thick,->]
    % Inertial Frame
    \coordinate (O)  at (0,0,0);
    \begin{scope}
      \draw (O)  -- +(1,0,0) node[          right= 0]{$\hat{e}_x$};
      \draw (O)  -- +(0,1,0) node[below=3,  right= 0]{$\hat{e}_y$};
      \draw (O)  -- +(0,0,1) node[below      left= 0]{$\hat{e}_z$};
    \end{scope}
    \node at (-0.5,0.5,-0.5) {Inertial Frame};
    % Non-Inertial Frame
    \coordinate (O2)  at (2,0,1);
    \begin{scope}[dotted,thin,-]
      \draw (O2) -- +(1,0,0) node[          right= 0]{}; %{$\hat{e}_x$};
      \draw (O2) -- +(0,1,0) node[below=3,  right= 0]{}; %{$\hat{e}_y$};
      \draw (O2) -- +(0,0,1) node[below      left= 0]{}; %{$\hat{e}_z$};
    \end{scope}
%
% Use the fortran code to rotate the unit vectors about the specified
% rotation vector.
% NOTE: The Fortran source to create this output is commented out at the
%       bottom of this file, and the output is 
%--------------------------------------------------------------------------
% OUTPUT FROM FORTRAN CODE
%--------------------------------------------------------------------------
% Info for the rotating frame:
%   angle_deg =        10.00000
%   angle_rad =         0.17453
%   axis      =         0.88465       0.44233       0.14744
%   omega     =         0.15440       0.07720       0.02573
%   ex        =         0.99670       0.03155      -0.07483
%   ey        =        -0.01966       0.98778       0.15461
%   ez        =         0.07879      -0.15263       0.98514
%--------------------------------------------------------------------------
    % Rotation Vector (not a unit vector)
    \begin{scope}[violet,thick,->]
      \draw (O2)  -- +(0.88,0.44,0.15) node[right= 0]{$\Omega$};
    \end{scope}
    % Rotated Unit Vectors
    \begin{scope}[red,thick,->]
%     \draw (O)  -- +( 0.99, 0.03,-0.07) node[          right= 0]{$\hat{e}_x$};
%     \draw (O)  -- +(-0.02, 0.99, 0.15) node[below=3,  right= 0]{$\hat{e}_y$};
%     \draw (O)  -- +( 0.08,-0.15, 0.99) node[below      left= 0]{$\hat{e}_z$};
      \draw (O2) -- +( 0.99, 0.03,-0.07) node[          right= 0]{$\hat{e}_{\hat{x}}$};
      \draw (O2) -- +(-0.02, 0.99, 0.15) node[below=3,  right= 0]{$\hat{e}_{\hat{y}}$};
      \draw (O2) -- +( 0.08,-0.15, 0.99) node[below      left= 0]{$\hat{e}_{\hat{z}}$};
    \end{scope}
%
    % Vectors
    \coordinate (X)  at (2.25,0.25,1.50);
    \begin{scope}[blue,thick,->]
      \draw (O)  -- (O2) node[midway, above=1, right=12]{${\mathbf{o}}$};
      \draw (O2) -- (X)  node[midway,          right   ]{$\hat{\mathbf{x}}$};
      \draw (O)  -- (X)  node[midway, above            ]{${\mathbf{x}}$};
    \end{scope}
    \node at (3.00,0.5,0.70) {Non-Inertial Frame};
  \end{scope}
\end{tikzpicture}

\end{document}

%###############################################################################
%###############################################################################
%###############################################################################

%!==========================================================================
%!> Program to compute rotated non-inertial reference frame for tikz image.
%!==========================================================================
%program main
%
%  implicit none
%
%  real, parameter      :: rad_per_deg = 3.14159 / 180.
%  real, parameter      :: deg_per_rad = 1. / rad_per_deg
%
%! rotation angle and axis
%  real                 :: rot_angle_deg = 10.0
%  real                 :: rot_angle_rad
%  real, dimension(3)   :: rot_axis = (/ 0.60, 0.30, 0.10 /)
%  real, dimension(3)   :: rot_vec
%  real, dimension(3,3) :: rot_ten
%
%! coordinate transformation tensors
%  real, dimension(3,3) :: Q, Q_inv
%
%! x,y,z unit vectors
%  real, dimension(3)   :: ex  = (/ 1.00, 0.00, 0.00 /)
%  real, dimension(3)   :: ey  = (/ 0.00, 1.00, 0.00 /)
%  real, dimension(3)   :: ez  = (/ 0.00, 0.00, 1.00 /)
%
%  continue
%
%! Rotation angle, (unit) axis, vector, and tensor.
%  rot_angle_rad = rot_angle_deg * rad_per_deg
%  rot_axis(:)   = rot_axis(:) / &
%                  sqrt( rot_axis(1)**2 + rot_axis(2)**2 + rot_axis(3)**2 )
%  rot_vec(:)    = rot_axis(:) * rot_angle_rad
%  rot_ten       = get_rot_ten( rot_vec )
%
%! Moving frame coordinate transformation tensor, $Q_{ij}$
%  Q     = get_rotation_matrix(rot_axis,rot_angle_rad) ! fix_to_rot
%  Q_inv = transpose(Q)                                ! rot_to_fix
%
%! Transform x,y,z unit vectors
%  ex = matmul( Q, ex )
%  ey = matmul( Q, ey )
%  ez = matmul( Q, ez )
%
%! Output
%  101 format(a,x,99(f14.5))
%  write(*,'(a)')   'Info for the rotating frame:'
%  write(*,101)     '  angle_deg = ', rot_angle_deg
%  write(*,101)     '  angle_rad = ', rot_angle_rad
%  call output_vec( '  axis      = ', rot_axis )
%  call output_vec( '  omega     = ', rot_vec )
%  call output_vec( '  ex        = ', ex )
%  call output_vec( '  ey        = ', ey )
%  call output_vec( '  ez        = ', ez )
%
%contains
%
%!==========================================================================
%!> Compute the moving frame rotation rate tensor from the rotation rate
%!> vector.
%function get_rot_ten( rot_vec ) &
%  result(rot_ten)
%
%  implicit none
%
%  real, dimension(3,3)           :: rot_ten
%  real, dimension(3), intent(in) :: rot_vec
%  
%  real, dimension(3,3,3) :: eijk
%  integer :: i,j,k
%
%! Permutation tensor, $\epsilon_{ijk}$
%  eijk(:,:,:) =  0.0
%  eijk(1,2,3) =  1.0
%  eijk(2,3,1) =  1.0
%  eijk(3,1,2) =  1.0
%  eijk(3,2,1) = -1.0
%  eijk(2,1,3) = -1.0
%  eijk(1,3,2) = -1.0
%
%! Moving frame rotation rate tensor, $\Omega_{ij}$
%  do k = 1, 3
%  do j = 1, 3
%  do i = 1, 3
%    rot_ten(i,j) = -eijk(i,j,k) * rot_vec(k)
%  enddo
%  enddo
%  enddo
%
%  return
%end function get_rot_ten
%!==========================================================================
%!
%!> Compute the transformation matrix to convert from Cartesian frame to
%!> a rotating frame defined by unit rotation axis and rotation angle.
%!>
%!> Ref: https://en.wikipedia.org/wiki/Rotation_matrix
%!
%function get_rotation_matrix(axis,angle) &
%  result (rot_mat)
%  !
%  !.. Formal Arguments ..
%  real, dimension(3,3)             :: rot_mat !< transformation matrix
%  real, dimension(3)  , intent(in) :: axis    !< unit vector of rotation axis
%  real                , intent(in) :: angle   !< angle in radians
%  !
%  !.. Local Scalars ..
%  real :: x, y, z, c, s, omc
%  !
%  ! Make axis a unit vector, if not already.
%  s = sqrt( axis(1)**2 + axis(2)**2 + axis(3)**2 )
%  x = axis(1) / s
%  y = axis(2) / s
%  z = axis(3) / s
%  c = cos(angle)
%  s = sin(angle)
%  omc = 1.0 - c
%  rot_mat = transpose( &
%              reshape( (/ x*x*omc +   c , x*y*omc - z*s, x*z*omc + y*s ,  &
%                          x*y*omc + z*s , y*y*omc +   c, y*z*omc - x*s ,  &
%                          x*z*omc - y*s , y*z*omc + x*s, z*z*omc +   c /) &
%                   , (/ 3, 3 /) ) &
%            )
%end function get_rotation_matrix
%!==========================================================================
%subroutine output_vec(text,v)
%  character*(*),      intent(in) :: text
%  real, dimension(3), intent(in) :: v
%  write(*,101) text,v
%  101 format(a,x,99(f14.5))
%end subroutine output_vec
%!==========================================================================
%end program main

